\documentclass[addpoints]{exam}
% \documentclass[addpoints,12pt]{exam}

\usepackage[slovene]{babel}
\usepackage{amsfonts}
\usepackage{amssymb}
\usepackage{amsmath}
\usepackage{float}
\usepackage{graphicx}
\usepackage{enumitem}
\usepackage{color}
\usepackage[gen]{eurosym}
\usepackage{newunicodechar}
\newunicodechar{€}{\euro}


%Komentar naslednje vrstice glede na verzijo testa -- rešitve ali prostor za odgovore:
% \printanswers

% \linespread{2}


\pointsinrightmargin
\marginbonuspointname{}


\renewcommand{\solutiontitle}{\noindent\textbf{Rešitve in točkovanje:}\par\noindent}

\colorgrids
\definecolor{GridColor}{gray}{0.7}
\setlength{\gridsize}{7mm}

\extrawidth{5mm}
\setlength{\rightpointsmargin}{1.5cm}

\begin{document}

\runningfooter{}{}{}

\begin{center}
    \fbox{\fbox{\parbox{15cm}{\centering
    Četrto pisno ocenjevanje znanja \\ 1. letnik -- test A \\ 1. april 2026 \\ (Realna števila)}}}
\end{center}

\vspace{0.1in}
\makebox[\textwidth]{Ime in priimek, oddelek:\enspace\hrulefill}


\begin{center}
    \chqword{Naloga}
    \chpword{Možne točke}
    \chbpword{Dodatne točke}
    \chsword{Dosežene točke}
    \chtword{Skupaj}  
    \combinedpointtable[h][questions]
    % \combinedgradetable[h][questions]

\end{center}

\vspace{0.1in}


\begin{questions}

    % 1. vprašanje

    \question[6]
    Izračunajte natančno vrednost izraza.
    $$ \sqrt{\sqrt{1024}}-\dfrac{3-\sqrt{2}}{\sqrt{2}-1}+\sqrt[3]{-8}+\left(2-\sqrt{2}\right)^2 $$

    \begin{solution}[7cm]
        \begin{align*} 
            \dfrac{\sqrt{5}}{\sqrt{5}+1}-\dfrac{5-3\sqrt{5}}{4\sqrt{5}}+\left(3+\sqrt{5}\right)^2-\sqrt{125}&=\dfrac{\sqrt{5}(\sqrt{5}-1)}{(\sqrt{5}+1)(\sqrt{5}-1)}-\dfrac{(5-3\sqrt{5})\sqrt{5}}{4\sqrt{5}\sqrt{5}}+9+6\sqrt{5}+5-5\sqrt{5}= \\
                                                                                                            &=\dfrac{5-\sqrt{5}}{5-1}-\dfrac{5\sqrt{5}-3\cdot 5}{4\cdot 5}+14+\sqrt{5}= \\
                                                                                                            &=\frac{5-\sqrt{5}}{4}-\dfrac{\sqrt{5}-3}{4}+14+\sqrt{5}= \\
                                                                                                            &=\frac{8-2\sqrt{5}}{4}+14+\sqrt{5}= \\
                                                                                                            &=16+\frac{1}{2}\sqrt{5}
        \end{align*}
        Za racionalizacijo $2$ točki, za izračun kvadrata dvočlenika $1$ točka, delno korenjenje $1$ točka, izračun $1$ točka, končne rezultat $1$ točka.
    \end{solution}
    
    
    % \newpage

    % 2. vprašanje
    \question[4]
    Rešite enačbo.
    $$ \dfrac{x-4}{x-2}+\dfrac{x+4}{x+1}=\dfrac{2x^2}{x^2-x-2} $$
    
    
    \begin{solution}[4cm]
        \begin{align*} 
            \dfrac{2}{x-3}-\dfrac{x}{x-1} &=\dfrac{2x}{x^2-4x+3} \\
            2(x-1)-x(x-3) &= 2x \\
            2x-2-x^2+3x-2x &=0 \\
            x^2-3x+2 &=0 \\
            (x-1)(x-2) &=0
        \end{align*}
        Pravilno odpravljanje ulomkov oziroma skupni imenovalec $1$ točka, zapis in upoštevanje pogojev za rešitev ($x\notin\{1,3\}$) $1$ točka, pravilno reševanje razcepne enačbe $1$ točka, zapis rešitve $x=2$ $1$ točka.
    \end{solution}

        

    \newpage
    % 3. vprašanje
    \question[5]
    Rešite sistem neenačb, rešitev grafično upodobite in zapišite z intervalom.
    $$ \left(3x-5<x+3\right)\lor\left(2x\geq x+6\right) $$
    
    
    \begin{solution}[12cm]
        Rešitev prve enačbe $x>2$ ali $x\in(2,\infty)$ $1$ točka, rešitev druge enačbe $x\geq 3$ ali $x\in[3,\infty)$ $1$ točka, rešitev sistema $x\geq 3$ ali $x\in[3,\infty)$ $1$ točka, zapis končne rešitve z intervalom $1$ točka, grafična upodobitev $1$ točka.
    \end{solution}


    


    % 4. vprašanje
    \question[4]
    Rešite sistem enačb.
    $$ \begin{aligned}
        y-x&=3 \\ x-y&=5
    \end{aligned} $$
    
    
    \begin{solution}[7cm]
        Pravilno reševanje sistema enačb do $2$ točki, pravilno sklepanje o neobstoju rešitve in zapis tega do $2$ točki.
    \end{solution}



    % 5. vprašanje
    % \question[4]
    % Prva cev sama napolni bazen v $6$ urah, druga sama pa v $9$ urah. Drugo cev odpremo eno uro kasneje kot prvo.
    % V kolikšnem času bo bazen poln?
    
    
    % \begin{solution}[8cm]
    %     ???
    % \end{solution}




    \newpage 

    
    % 6. vprašanje
    \question[4]
    Natančno izračunajte.
    $$ \left\lvert \lvert\sqrt{2}+4\rvert-\sqrt{2}\right\rvert + \left\lvert \lvert\sqrt{2}-4\rvert-\sqrt{2}\right\rvert $$
    
    
    
    \begin{solution}[11cm]
        \begin{align*}
            \left\lvert \lvert\sqrt{2}+4\rvert-\sqrt{2}\right\rvert + \left\lvert \lvert\sqrt{2}-4\rvert-\sqrt{2}\right\rvert &=
            \left\lvert \sqrt{2}+4-\sqrt{2}\right\rvert + \left\lvert -\sqrt{2}+4-\sqrt{2}\right\rvert= \\
            &= \left\lvert 4\right\rvert + \left\lvert 4-2\sqrt{2}\right\rvert= \\
            &= 4+4-2\sqrt{2}= \\
            &= 8-2\sqrt{2}
        \end{align*}
        Pravilno razreševanje absolutnih vrednosti do $3$ točke, končen rezultat $1$ točka.
    \end{solution}




    % 7. vprašanje
    \question[5]
    Rešite enačbo in opišite njen geometrijski pomen ter rešitev grafično prikažite.
    $$  \left\lvert x-2\right\rvert=3 $$
    
    
    \begin{solution}[7cm]
        Opis geometrijskega pomena (oddaljenost števila $x$ od števila $-3$ je $2$) $1$ točka, grafična predstavitev $1$ točka, reševanje enačbe z upoštevanjem pogojev do $2$ točke, zapis rešitve $x\in\{-5,-1\}$ $1$ točka..
    \end{solution}


    \newpage

    % \makebox[\textwidth]{Ime in priimek, oddelek:\enspace\hrulefill}

    % ~\\



    % 8. vprašanje
    % \question[5]
    % Izračunajte vrednosti izraza glede na vrednosti spremenljivke $x$.
    % $$\left\lvert x+8 \right\rvert + \left\lvert x-2 \right\rvert $$
    
    
    % \begin{solution}[10cm]
    %     ???
    % \end{solution}



    % 9. vprašanje
    \question[5]
    Prvi delavec sam pozida steno v $10$ urah, drugi v $12$ urah, tretji v $8$ urah. Delavci
    skupaj začnejo zidati steno. Po dveh urah tretji delavec odide, pridruži pa se četrti delavec.
    Skupaj s prvim in drugim delavcem nato končajo steno v eni uri. V kolikšnem času četrti delavec
    pozida steno?
    
    
    \begin{solution}[13cm]
        Zapis enačbe po besedilu $1$ točka, pravilno reševanje enačbe $1$ točka, zapis pravilne rešitve $1$ točka, zapis pravilnega odgovora $1$ točka.
    \end{solution}






    
    %  \newpage



    % 3. vprašanje
    % \question[4]
    % Produkt dveh števil je $192$, njun največji skupni delitelj pa $4$. Poiščite števili.
    
    % \begin{solution}[7cm]
    %     \begin{align*}
    %         1-(1+x^8)(1+x^4)(1+x^2)(1+x)(1-x)&=1-(1+x^8)(1+x^4)(1+x^2)(1-x^2)= \\
    %             &=1-(1+x^8)(1+x^4)(1-x^4)= \\
    %             &=1-(1+x^8)(1-x^8)= \\
    %             &=1-(1-x^{16})= \\
    %             &=1-1+x^{16}= \\
    %             &=x^{16}
    %     \end{align*}
    %     Za pravilno uporabo pravila razlike kvadratov $1$ točka, za pravilno upoštevan negativen predznak $1$ točka, končen rezultat $1$ točka.
    % \end{solution}


    % \begin{description}
    %     \item[a] $(x-a)^2+(y-b)^2-c^2=0$ 
    %     \item[b] $d^2(x-a)^2+f^2(y-b)^2-(df)^2=0$
    %     \item[c] $S(a,b)$
    %     \item[d] $e^2=d^2-f^2;\ d>f$
    % \end{description}
    % Krožnica: \fillin[a c][2cm]
    % \\ Elipsa:  \fillin[b c d][2cm]




    %  \newpage

    %  \noindent Ime in priimek: \underline{~~~~~~~~~~~~~~~~~~~~~~~~~~~~~~~~~~~~~~~~~~~~~~~~}

    %  ~\\

      

    % dodatno vprašanje
    \bonusquestion[3]
    \textit{Dodatna naloga:} Natančno izračunajte. $$ \left(\sqrt{3}+\sqrt{5}\right)\sqrt{8-2\sqrt{15}} $$

    \begin{solution}[5cm]
        Pravilen izračun do $2$ točki.
    \end{solution}




\end{questions}



\end{document}